\section{Introduction}
\label{sec:introduction}
Modern text-to-speech (TTS) and voice cloning produce speech hard to distinguish
from genuine recordings, raising risks of fraud, impersonation and
misinformation~\cite{yi2023} and motivating watermarks that trace audio to its
source~\cite{san2024,chen2023}. We state up front which part of that problem we
address. \texttt{RAWMER} is \emph{provider-side}: only the generating provider,
holding the secret key and a per-utterance reference, can verify, and only for audio
it generated. It substantiates or refutes an ownership claim about a specific
candidate generation record; it cannot flag a clone from a non-cooperating system,
which remains the province of blind detection~\cite{san2024} and deepfake
classifiers~\cite{yi2023}---the two are complementary
(\S\ref{sec:method:threat}).

\textbf{Related work.} Post-hoc waveform watermarks such as
AudioSeal~\cite{san2024}, WavMark~\cite{chen2023} and SilentCipher~\cite{singh2024}
embed after synthesis and decode blindly. In-generation schemes tie the mark to the pipeline instead: Timbre
Watermarking~\cite{liu2024timbre} into speaker timbre,
Trace\-able\-Speech~\cite{zhou2024} into discrete speech tokens,
Groot~\cite{liu2024groot} into a diffusion generator's noise prior. Since most TTS
systems emit a mel spectrogram before a vocoder~\cite{kong2020,kong2021}, that mel is
a shared, model-agnostic interface needing no vocoder retraining.
MelShield~\cite{jin2026} follows this route, embedding a keyed payload via low-energy
spread-spectrum perturbations and verifying against a retained clean reference. We
adopt the same setting and take MelShield---concurrent work, re-implemented
locally---as our primary baseline.

Our point of departure is the encoding mechanism. MelShield~\cite{jin2026} and prior
mel-domain schemes represent each bit as an absolute perturbation pattern, detected by
correlating the reference residual with the keyed pattern. Absolute residuals are
sensitive to gain, frequency shaping and vocoder reconstruction error; at lower SNRs
the pattern is buried and bit accuracy falls toward chance. This fragility under
strong additive noise---not under compression or filtering---is our target.

Our remedy rests on a principle that is not new. Encoding a bit as the sign of an
energy difference between two keyed sets of host samples is the classical
\emph{patchwork} construction~\cite{bender1996,yeo2003}, one of a lineage of
relational and quantization-based embeddings including QIM~\cite{chen2001}; our
detector is likewise an \emph{informed} one~\cite{cox2008}, since it reads the host.
Our claim is narrower, and has two parts.
\textbf{Relative mel-energy encoding before vocoding:} moving patchwork-style
differential encoding into the pre-vocoder mel domain yields a score $z_i$ in which
common-mode distortion cancels; alone this already exceeds---narrowly---a
quality-matched absolute-residual baseline under strong noise.
\textbf{Reliability-aware bin-group selection:} the clean reference, already required
here, picks among $C$ keyed candidates the pair whose groups both have dynamic-range
headroom and stable energy and are balanced in reliability (Eq.~\eqref{eq:select});
this is the larger contributor. On HiFi-GAN~\cite{kong2020} and
DiffWave~\cite{kong2021} the combination improves noise robustness over a
quality-matched MelShield~\cite{jin2026} re-implementation, with chance-level false
positives, $\approx$\,22~KB references and verification from unaligned 25\% clips.
It does \emph{not} withstand an attacker who desynchronizes time or frequency.
